\documentclass[a4paper]{article}
\usepackage[pdftex]{graphicx}
\usepackage[utf8]{inputenc}
\usepackage{enumerate}
\usepackage{amssymb}
\usepackage{colortbl}
\usepackage{icomma}
\usepackage{colortbl}
\usepackage{musixtex}
\usepackage{siunitx}
\sisetup{locale=DE} 
\usepackage{geometry}
\geometry{a4paper, top=15mm, left=15mm, right=15mm, bottom=15mm,
	headsep=10mm, footskip=12mm}
\usepackage{href-ul}
\hypersetup{
	colorlinks=true,
	linkcolor=blue,
	urlcolor=blue}
\begin{document}

{\bf Die Tonleiter, mathematisch gesehen}
\begin{enumerate}[1.]
{\bf \item Fülle die Tabelle mit Hilfe der Formel $y =440 \cdot 2^{\frac{x}{12}}$ }  aus.
 Hierbei ist x die Anzahl der Halbtonschritte und y die Frequenz.

\renewcommand{\arraystretch}{2}
\begin{tabular}{|>{\columncolor[gray]{.8}}p{1.5 cm}|p{1.5 cm}|p{1.5 cm}|p{1.5 cm}|p{1.5 cm}|p{1.5 cm}|p{1.5 cm}|p{1.5 cm}|}
	\hline
	\rowcolor[gray]{.8}	 Ton &  a, & c & c\# & d & d\# & e & f\\
	\hline
	 x & -12  & -9  &  -8 &-7 & -6& -5 &  -4 \\
	\hline 
	 y in Hz &  & & & & & & \\ 
	 \hline
\end{tabular}


\renewcommand{\arraystretch}{2}
\begin{tabular}{|>{\columncolor[gray]{.8}}p{1.5 cm}|p{1.5 cm}|p{1.5 cm}|p{1.5 cm}|p{1.5 cm}|p{1.5 cm}|p{1.5 cm}|p{1.5 cm}|p{1.5 cm}|}
	\hline
	\rowcolor[gray]{.8}	 Ton &  f\#  & g & g\# & a & a\# & h & c' & a'\\
	\hline
	x & -3  & -2 &  -1 & 0 & 1 & 2 &  3 & 12 \\
	\hline 
	y in Hz &  & & & 440 & & & &\\ 
	\hline
\end{tabular}


{\bf \item  Stelle die Frequenzen über den Halbtonschritten} grafisch dar:\\
a) mit linearer y-Achse\\
\begin{center}
\hspace{-0.4 cm}
\includegraphics[width=15 cm]{expto01.png}	\\


\begin{music}
	\parindent10mm
	\instrumentnumber{1}
	\setstaffs1{1}
	\startextract
	\Notes 

	\qa {a}
	\hskip 1  cm
	\qa c
	\qa{^c}	
	\qa d  
	\qa{^d}
	\qa e 
	\qa f 
	\qa{^f}
	\qa {g'}
	\qa{^a}	
	\qa{a}
	\qa b 
	\qa {c'}
	\hskip 2 cm
	\qa{a''} 
	\en
	\endextract
\end{music}

b) mit logarithmischer y-Achse\\

\hspace{-0.4 cm}
\includegraphics[width=15 cm]{logto01.png}	\\
\end{center}

\begin{music}
	\parindent10mm
	\instrumentnumber{1}
	\setstaffs1{1}
	\startextract
	\Notes 
	\qa c
	\qa{^c}	
	\qa d  
	\qa{^d}
	\qa e 
	\qa f 
	\qa{^f}
	\qa g 
	\qa{^g}
	\qa {'a} 
	\qa{^a}
	\qa b 
	\qa {c'}
	\en
	\endextract
\end{music}
{\bf \item Fülle die Tabelle mit den musikalischen Intervallen aus}

\renewcommand{\arraystretch}{1}
\begin{tabular}{|>{\columncolor[gray]{.8}}p{3 cm}|p{3 cm}|p{3 cm}|p{3 cm}|p{3 cm}|}
	\hline
	\rowcolor[gray]{.8}	Name & Halbtonschritte &  Noten & $2^{x/12}$ & Frequenzverhältnis   \\
	\hline
     Prime & & 	 
	  \begin{music}
	  	\instrumentnumber{1}
	  	\setstaffs1{1}
	  	\startextract
	  	\Notes \zcharnote{c}{\raisebox{-3ex}{C}} \qa c \hskip 1.2cm  \en
	  	\endextract
	  \end{music}
	  &  &\\
	\hline
	  kleine Sekunde & & 
	  \begin{music}
	  	\instrumentnumber{1}
	  	\setstaffs1{1}
	  	\startextract
	  	\Notes \zcharnote{c}{\raisebox{-3ex}{C}} \qa c \hskip 1.2cm  \en
	  	\endextract
	  \end{music}
	  & & \\
	\hline
	 große Sekunde & & 
	 \begin{music}
	 	\instrumentnumber{1}
	 	\setstaffs1{1}
	 	\startextract
	 	\Notes \zcharnote{c}{\raisebox{-3ex}{C}} \qa c \hskip 1.2cm  \en
	 	\endextract
	 \end{music}
	 & & \\
	\hline
	  kleine Terz & & 
	  \begin{music}
	  	\instrumentnumber{1}
	  	\setstaffs1{1}
	  	\startextract
	  	\Notes \zcharnote{c}{\raisebox{-3ex}{C}} \qa c \hskip 1.2cm  \en
	  	\endextract
	  \end{music}
	  & & \\
	\hline
	 große Terz & &
	 \begin{music}
	 	\instrumentnumber{1}
	 	\setstaffs1{1}
	 	\startextract
	 	\Notes \zcharnote{c}{\raisebox{-3ex}{C}} \qa c \hskip 1.2cm  \en
	 	\endextract
	 \end{music}
	  & & \\
	\hline
	 Quarte & & 
	 \begin{music}
	 	\instrumentnumber{1}
	 	\setstaffs1{1}
	 	\startextract
	 	\Notes \zcharnote{c}{\raisebox{-3ex}{C}} \qa c \hskip 1.2cm  \en
	 	\endextract
	 \end{music}
	 & & \\
	\hline
\end{tabular}

\vspace{1 cm}
Frequenzverhältnisse:
\quad\fbox{1 : 1}
\quad\fbox{5:4}
\quad\fbox{4:3}
\quad\fbox{6:5}
\quad\fbox{9:8}
\quad\fbox{16:15}

\newpage
\begin{music}
	\parindent10mm
	\instrumentnumber{1}
	\setstaffs1{1}
	\startextract
	\Notes 
	\qa c
	\qa{^c}	
	\qa d  
	\qa{^d}
	\qa e 
	\qa f 
	\qa{^f}
	\qa g 
	\qa{^g}
	\qa {'a} 
	\qa{^a}
	\qa b 
	\qa {c'}
	\en
	\endextract
\end{music}
\renewcommand{\arraystretch}{1}
\begin{tabular}{|>{\columncolor[gray]{.8}}p{3 cm}|p{3 cm}|p{3 cm}|p{3 cm}|p{3 cm}|}
	\hline
	\rowcolor[gray]{.8}	Name & Halbtonschritte &  Noten & $2^{x/12}$ & Frequenzverhältnis  \\
	\hline
	  Tritonus & & 
	  \begin{music}
	  	\instrumentnumber{1}
	  	\setstaffs1{1}
	  	\startextract
	  	\Notes \zcharnote{c}{\raisebox{-3ex}{C}} \qa c \hskip 1.2cm  \en
	  	\endextract
	  \end{music}
	  & & \\
	\hline
	  Quinte & & 
	  \begin{music}
	  	\instrumentnumber{1}
	  	\setstaffs1{1}
	  	\startextract
	  	\Notes \zcharnote{c}{\raisebox{-3ex}{C}} \qa c \hskip 1.2cm  \en
	  	\endextract
	  \end{music}
	  &  & \\
	\hline
	 kleine Sexte & & 
	 \begin{music}
	 	\instrumentnumber{1}
	 	\setstaffs1{1}
	 	\startextract
	 	\Notes \zcharnote{c}{\raisebox{-3ex}{C}} \qa c \hskip 1.2cm  \en
	 	\endextract
	 \end{music}
	 & & \\
	\hline
	  große Sexte & & 
	  \begin{music}
	  	\instrumentnumber{1}
	  	\setstaffs1{1}
	  	\startextract
	  	\Notes \zcharnote{c}{\raisebox{-3ex}{C}} \qa c \hskip 1.2cm  \en
	  	\endextract
	  \end{music}
	  & & \\
	\hline
	 kleine Septime & &
	 \begin{music}
	 	\instrumentnumber{1}
	 	\setstaffs1{1}
	 	\startextract
	 	\Notes \zcharnote{c}{\raisebox{-3ex}{C}} \qa c \hskip 1.2cm  \en
	 	\endextract
	 \end{music}
	  & & \\
	\hline
	  große Septime & & 
	  \begin{music}
	  	\instrumentnumber{1}
	  	\setstaffs1{1}
	  	\startextract
	  	\Notes \zcharnote{c}{\raisebox{-3ex}{C}} \qa c \hskip 1.2cm  \en
	  	\endextract
	  \end{music}
	  & & \\
	\hline
	  Oktave & & 
	  \begin{music}
	  	\instrumentnumber{1}
	  	\setstaffs1{1}
	  	\startextract
	  	\Notes \zcharnote{c}{\raisebox{-3ex}{C}} \qa c \hskip 1.2cm  \en
	  	\endextract
	  \end{music}
	  & & \\
	\hline
\end{tabular}

\vspace{1 cm}
 Frequenzverhältnisse:
	\quad \fbox{3:2}
\quad \fbox{5:3}
\quad \fbox{45:32}
\quad \fbox{2:1}
\quad \fbox{9:5}
\quad \fbox{8:5}
\quad \fbox{15:8}

\vspace{1 cm}
{\bf \item Um wieviel Prozent steigt die Frequenz bei \underbar{einem} Halbtonschritt?}
\vspace{1 cm}
{\bf \item Spiele die Intervalle auf einer Piano-App}.


\end{enumerate} 

\fbox{
\begin{minipage}{0.5\textwidth}
	Zur Lösung bitte \href{https://www.okuyakl.de/math/m10tonL135/ll135.pdf}{hier klicken} oder den QR-Code scannen.\\
	Weitere Arbeitsblätter gibt es unter 
	
	\href{https://www.okuyakl.de}{www.okuyakl.de}
\end{minipage}
\hfill
\begin{minipage}{0.4\textwidth}
	\includegraphics[width=1.5 cm]{../../viecher/zwe03}
	\includegraphics[width=3 cm]{qr135}
	\includegraphics[width=2 cm]{../../viecher/afanticon1}
	
	\end{minipage}}
	
\end{document}